\chapter{Your First CNA Game}
\label{ch:first-game}

This chapter walks through the smallest complete CNA program --- adapted from the minimal
XNA-style skeleton in CNA's own README, with the base-loop calls made explicit --- line by line,
as a concrete anchor for the framework
concepts explained more fully in the chapters that follow.

\section{The whole program}

\begin{lstlisting}
#include <memory>

#include "Microsoft/Xna/Framework/Game.hpp"
#include "Microsoft/Xna/Framework/Color.hpp"
#include "Microsoft/Xna/Framework/Graphics/GraphicsDeviceManager.hpp"
#include "Microsoft/Xna/Framework/Graphics/SpriteBatch.hpp"
#include "Microsoft/Xna/Framework/Graphics/Texture2D.hpp"

using namespace Microsoft::Xna::Framework;
using namespace Microsoft::Xna::Framework::Graphics;

class MyGame final : public Game {
public:
    MyGame()
        : graphics_(this)
    {
    }

protected:
    void LoadContent() override
    {
        spriteBatch_ = std::make_unique<SpriteBatch>(getGraphicsDeviceProperty());
        logo_ = std::make_unique<Texture2D>("assets/logo.png", getGraphicsDeviceProperty());
    }

    void Update(GameTime& gameTime) override
    {
        // Update game state here.
        Game::Update(gameTime);
    }

    void Draw(const GameTime& gameTime) override
    {
        auto& device = getGraphicsDeviceProperty();
        device.Clear(CornflowerBlue);

        spriteBatch_->Begin();
        spriteBatch_->Draw(*logo_, 100.0f, 80.0f);
        spriteBatch_->End();
        Game::Draw(gameTime);
    }

private:
    GraphicsDeviceManager graphics_;
    std::unique_ptr<SpriteBatch> spriteBatch_;
    std::unique_ptr<Texture2D> logo_;
};

int main()
{
    MyGame game;
    game.Run();
    return 0;
}
\end{lstlisting}

\section{Reading it top to bottom}

\textbf{The includes are the API surface, and nothing else.} Every header included is a real
XNA~4.0 header, under the real XNA namespace path --- there is no CNA-specific include
required just to boot a game. This is the API-mirroring goal from
Chapter~\ref{ch:design-philosophy} made visible: a developer who has written XNA or FNA code
before can read this file without learning a single CNA-specific concept first.

\textbf{\texttt{MyGame} inherits \cnaclass{Game}, and owns a
\cnaclass{GraphicsDeviceManager} constructed against itself.} This is the standard XNA
source-level pattern, but CNA's C++ ownership differs from FNA's: the \cnaclass{Game} base
already default-constructs its own \cnaclass{GraphicsDevice} before the derived constructor
creates this manager. The manager registers the lifecycle services and, during
\texttt{Game::DoInitialize()}, applies the preferred settings by resetting that same game-owned
object; it does not construct or own a second device. Chapter~\ref{ch:game-loop} covers that
boot sequence in full, and Chapter~\ref{ch:graphicsdevice} covers what the resulting
\cnaclass{GraphicsDevice} exposes.

\textbf{\texttt{LoadContent()} is where GPU-backed resources are created, not the
constructor.} By the time \texttt{LoadContent()} runs, the graphics renderer has already been
constructed by the base \cnaclass{Game} and then configured by
\cnaclass{GraphicsDeviceManager}, so it is safe to construct a \cnaclass{SpriteBatch} against
\texttt{getGraphicsDeviceProperty()} and to load a \cnaclass{Texture2D} from a file path. A
derived constructor can technically already see CNA's default renderer, but it precedes the
manager's preferred-setting reset and is the wrong lifecycle point for content. Defer resources
to \texttt{LoadContent()} for the stable, configured device rather than relying on the old,
incorrect claim that no renderer exists yet. This ordering constraint --- and CNA's
raw-asset-plus-JSON-descriptor content model that makes
\texttt{"assets/logo.png"} a literal, directly loadable file path rather than a compiled
\texttt{.xnb} content-pipeline output --- is covered fully in
Chapter~\ref{ch:content-and-assets}.

\textbf{\texttt{Update(GameTime\&)} and \texttt{Draw(const GameTime\&)} are the two halves of
every frame, and note both the constness difference and the base calls.} \texttt{Update} takes a mutable
\texttt{GameTime\&} (some XNA game loops advance simulation-only clocks here);
\texttt{Draw} takes a \texttt{const GameTime\&}, since rendering should not mutate timing
state. Calling \texttt{Game::Update(gameTime)} runs registered updateable components and the
framework dispatcher; calling \texttt{Game::Draw(gameTime)} runs registered drawable components.
Omitting the first call silently stops dynamic audio, media, and touch polling. Chapter~\ref{ch:game-loop}
covers the full lifecycle and the ordering choices this pair of overrides plugs into.

\textbf{The draw call is the canonical \cnaclass{SpriteBatch} triad:
\texttt{Begin()} / \texttt{Draw(...)} / \texttt{End()}.} \cnaclass{GraphicsDevice::Clear} takes a
named XNA color constant (\cnaclass{Color::CornflowerBlue} --- traditionally XNA's default
clear color in generated project templates, preserved here rather than swapped for something
``more CNA''). Application \texttt{Draw()} does not call \cnaclass{GraphicsDevice::Present()}
itself: after the override returns, \texttt{Game::EndDraw()} delegates to the manager, which
presents exactly once regardless of which public renderer identity from
Part~\ref{part:renderers} is active.
Chapter~\ref{ch:spritebatch} covers the full \cnaclass{SpriteBatch} API, including the
overloads of \texttt{Draw} beyond the two-float-position form used here.

\textbf{\texttt{main()} is exactly two lines of game-specific code.} Construct the game,
call \texttt{Run()}. Everything about window creation, renderer initialization, the timing
loop, and shutdown lives inside \cnaclass{Game::Run()} --- covered in
Chapter~\ref{ch:game-loop} --- not in application code. This is deliberate: it is the same
shape a real XNA \texttt{Program.cs}'s \texttt{Main} method has, preserved so that porting an
existing XNA game's entry point is close to a mechanical translation.

\section{A second program: making it move}
\label{sec:first-game-movement}

The program above proves a texture can reach the screen; it says nothing about the
\texttt{Update} / \texttt{Draw} split actually \emph{doing} anything, since its own
\texttt{Update} is empty. This section extends exactly the same skeleton --- same includes,
same class shape, same \texttt{main()} --- with the smallest genuinely interactive addition:
moving the logo with the arrow keys, at a constant real-world speed regardless of frame rate.
Three small, real changes, each grounded in the actual headers rather than invented:

\begin{lstlisting}
#include "Microsoft/Xna/Framework/Input/Keyboard.hpp"
#include "Microsoft/Xna/Framework/Vector2.hpp"

using namespace Microsoft::Xna::Framework::Input;

class MyGame final : public Game {
    // ... constructor and LoadContent() unchanged ...

    void Update(GameTime& gameTime) override
    {
        const KeyboardState keys = Keyboard::GetState();
        const float seconds =
            static_cast<float>(gameTime.getElapsedGameTimeProperty().getTotalSecondsProperty());
        const float speed = 200.0f; // pixels per second

        if (keys.IsKeyDown(Keys::Left))  position_.X -= speed * seconds;
        if (keys.IsKeyDown(Keys::Right)) position_.X += speed * seconds;
        if (keys.IsKeyDown(Keys::Up))    position_.Y -= speed * seconds;
        if (keys.IsKeyDown(Keys::Down))  position_.Y += speed * seconds;

        Game::Update(gameTime);
    }

    void Draw(const GameTime& gameTime) override
    {
        auto& device = getGraphicsDeviceProperty();
        device.Clear(CornflowerBlue);

        spriteBatch_->Begin();
        spriteBatch_->Draw(*logo_, position_, Color::White);
        spriteBatch_->End();
        Game::Draw(gameTime);
    }

private:
    GraphicsDeviceManager graphics_;
    std::unique_ptr<SpriteBatch> spriteBatch_;
    std::unique_ptr<Texture2D> logo_;
    Vector2 position_{100.0f, 80.0f}; // new: replaces the two hardcoded floats
};
\end{lstlisting}

\textbf{Reading polled input is a static call, not an object you own.} \cnaclass{Keyboard} has
no constructor to call and nothing to store between frames --- \texttt{Keyboard::GetState()}
is a static method returning a fresh, immutable \cnaclass{KeyboardState} snapshot every time
it is called, and \texttt{Update} calls it once per frame. This is the same shape
Chapter~\ref{ch:input-system} covers for \cnaclass{Mouse} and \cnaclass{GamePad}.

\textbf{Frame-rate independence comes from \texttt{GameTime}, not a hardcoded per-frame
step.} Multiplying a pixels-per-second constant by
\texttt{gameTime.get\allowbreak{}ElapsedGameTimeProperty().get\allowbreak{}TotalSecondsProperty()} --- the real
\cnaclass{TimeSpan} accessor pair, not an invented shortcut --- means the logo moves at the
same real-world speed whether the game is running at 30 FPS or 240 FPS. A version that instead
moved by a flat constant every \texttt{Update} call would move four times faster on a display
with a four times higher refresh rate, which is precisely the class of bug frame-rate-independent
movement exists to avoid.

\begin{sourcenote}
\textbf{Why this program uses a different \texttt{Draw} overload than the first one.} The
first program in this chapter calls \texttt{spriteBatch\_->Draw(*logo\_, 100.0f, 80.0f)} ---
but that specific two-float-position overload, reading
\texttt{modules/graphics/include/\allowbreak{}Microsoft/\allowbreak{}Xna/\allowbreak{}Framework/\allowbreak{}Graphics/\allowbreak{}SpriteBatch.hpp} directly, is tagged
\texttt{CNAEXT}: a CNA-only convenience overload, not part of real XNA~4.0
(Chapter~\ref{ch:design-philosophy} covers what that tag means). This second program instead
calls \texttt{Draw(*logo\_, position\_, Color::White)}, the real, unmarked XNA~4.0 overload
that takes a \cnaclass{Vector2} position and an explicit tint color --- worth calling out
directly here, since a reader who wants code that ports cleanly to real XNA/FNA should reach
for this overload as the default, not the convenience one the first, minimal example used to
stay as short as possible.
\end{sourcenote}

\section{A third program: bouncing off the edges, with sound}

The second program moved a sprite; it never asked whether that sprite had actually left the
visible screen. This section adds exactly that check, plus a real audio cue when it happens
--- the classic ``first game'' milestone of something actually \emph{reacting} to the game
state, not just rendering it. Three more real, grounded additions on top of the same running
example:

\begin{lstlisting}
class MyGame final : public Game {
    // ... constructor unchanged ...

    void LoadContent() override
    {
        spriteBatch_ = std::make_unique<SpriteBatch>(getGraphicsDeviceProperty());
        logo_ = std::make_unique<Texture2D>("assets/logo.png", getGraphicsDeviceProperty());
        bounceSound_ = std::make_unique<SoundEffect>("assets/bounce.wav");
    }

    void Update(GameTime& gameTime) override
    {
        const KeyboardState keys = Keyboard::GetState();
        const float seconds =
            static_cast<float>(gameTime.getElapsedGameTimeProperty().getTotalSecondsProperty());
        const float speed = 200.0f;

        if (keys.IsKeyDown(Keys::Left))  position_.X -= speed * seconds;
        if (keys.IsKeyDown(Keys::Right)) position_.X += speed * seconds;
        if (keys.IsKeyDown(Keys::Up))    position_.Y -= speed * seconds;
        if (keys.IsKeyDown(Keys::Down))  position_.Y += speed * seconds;

        const Viewport& viewport = getGraphicsDeviceProperty().getViewportProperty();
        const float maxX = static_cast<float>(viewport.getWidthProperty())  - logo_->getWidthProperty();
        const float maxY = static_cast<float>(viewport.getHeightProperty()) - logo_->getHeightProperty();

        bool hitEdge = false;
        if (position_.X < 0.0f)   { position_.X = 0.0f; hitEdge = true; }
        if (position_.X > maxX)   { position_.X = maxX; hitEdge = true; }
        if (position_.Y < 0.0f)   { position_.Y = 0.0f; hitEdge = true; }
        if (position_.Y > maxY)   { position_.Y = maxY; hitEdge = true; }

        if (hitEdge) bounceSound_->Play();
        Game::Update(gameTime);
    }

private:
    // ... other members unchanged ...
    std::unique_ptr<SoundEffect> bounceSound_;
};
\end{lstlisting}

\textbf{The bounds check clamps position, it does not just detect it.} Setting
\texttt{position\_.X} back to \texttt{0.0f} or \texttt{maxX} in the same branch that detects
the violation matters: without it, holding a direction key against the edge would let
\texttt{position\_} keep accumulating past the boundary every frame, and the sprite would need
to travel back through the entire out-of-bounds distance before visibly re-entering the
screen. Clamping in the same statement that detects the crossing keeps the visible sprite
exactly at the edge, every frame, for as long as the key is held.

\textbf{The viewport, not a hardcoded screen size, is what \texttt{maxX} / \texttt{maxY} are
computed against.} \cnaclass{Viewport::Width} / \texttt{Height} (Chapter~\ref{ch:graphicsdevice})
reflect whatever the back buffer's actual current size is --- reading them here, rather than
hardcoding the numbers the game happened to launch with, is what keeps this bounds check
correct if the window is resized or the game runs under a different
\cnaclass{PresentationMode} (Chapter~\ref{ch:game-loop}).

\textbf{\cnaclass{SoundEffect}'s file-path constructor is a \texttt{CNAEXT} convenience, the
same shape as \cnaclass{Texture2D}'s.} Real XNA has no such constructor either --- a real XNA
game always loads a \cnaclass{SoundEffect} through the compiled content pipeline. \texttt{Play()}
here uses the simplest zero-argument overload (default volume, pitch, and pan); a real game
that wants to vary bounce loudness by impact speed would reach for the three-argument
\texttt{Play(volume, pitch, pan)} overload instead --- covered in full in this book's audio
chapter.

\section{A fourth program: a real collision check}

The third program reacted to the screen's own edges; this section adds the other kind of
reaction a first game almost always needs next --- two sprites reacting to \emph{each other}.
A second, static texture (a coin) sits on screen; when the moving logo's own bounding rectangle
overlaps it, the coin is collected and a counter increments. One more real, grounded addition
on top of the same running example:

\begin{lstlisting}
class MyGame final : public Game {
    // ... constructor unchanged ...

    void LoadContent() override
    {
        spriteBatch_ = std::make_unique<SpriteBatch>(getGraphicsDeviceProperty());
        logo_ = std::make_unique<Texture2D>("assets/logo.png", getGraphicsDeviceProperty());
        bounceSound_ = std::make_unique<SoundEffect>("assets/bounce.wav");
        coin_ = std::make_unique<Texture2D>("assets/coin.png", getGraphicsDeviceProperty());
    }

    void Update(GameTime& gameTime) override
    {
        // ... movement and edge-clamping from the third program, unchanged ...

        if (!collected_)
        {
            const Rectangle logoBounds(static_cast<int>(position_.X),
                                        static_cast<int>(position_.Y),
                                        logo_->getWidthProperty(),
                                        logo_->getHeightProperty());
            const Rectangle coinBounds(static_cast<int>(coinPosition_.X),
                                        static_cast<int>(coinPosition_.Y),
                                        coin_->getWidthProperty(),
                                        coin_->getHeightProperty());
            if (logoBounds.Intersects(coinBounds))
            {
                collected_ = true;
                ++score_;
                bounceSound_->Play();
            }
        }

        Game::Update(gameTime);
    }

    void Draw(const GameTime& gameTime) override
    {
        auto& device = getGraphicsDeviceProperty();
        device.Clear(CornflowerBlue);

        spriteBatch_->Begin();
        spriteBatch_->Draw(*logo_, position_, Color::White);
        if (!collected_) spriteBatch_->Draw(*coin_, coinPosition_, Color::White);
        spriteBatch_->End();
        Game::Draw(gameTime);
    }

private:
    // ... other members unchanged ...
    std::unique_ptr<Texture2D> coin_;
    Vector2 coinPosition_{300.0f, 200.0f};
    bool collected_ = false;
    int score_ = 0;
};
\end{lstlisting}

\textbf{The bounding rectangles are built fresh every frame, from the same position/size data
already driving movement and drawing --- not stored and kept in sync by hand.}
\cnaclass{Rectangle}'s constructor (Chapter~\ref{ch:math-core-types}) takes an integer
\texttt{x} / \texttt{y} / \texttt{width} / \texttt{height}, so \texttt{position\_}'s float
components need an explicit \texttt{static\_cast<int>} --- the same pattern
\cnaclass{Texture2D::Width} / \texttt{Height} already supply directly for the size half of each
rectangle. Building both rectangles inline in \texttt{Update()}, from the exact position
fields \texttt{Draw()} also reads, is what keeps the collision check and the visible sprites
from ever silently disagreeing about where anything actually is.

\begin{sourcenote}
\textbf{A real, test-verified edge case worth knowing before relying on
\texttt{Intersects()}: touching is not overlapping.} \cnaclass{Rectangle}'s own real test
suite (\texttt{RectangleTests.cpp}) checks this precisely: two $100\times100$ rectangles
placed at $(0,0)$ and $(100,0)$ --- sharing an edge exactly, with zero pixels of true overlap
--- return \texttt{false} from \texttt{Intersects()}, the same as two rectangles with real
space between them. Only genuine, non-zero-area overlap counts. A first game that places a
coin exactly at the logo's own width, expecting the two to register as ``touching'' the
instant they meet, will see the check silently never fire; leaving at least a few pixels of
deliberate overlap margin between two sprites meant to collide avoids this class of bug
entirely.
\end{sourcenote}

\section{What happens if you change the renderer}

Nothing in the file above mentions a graphics renderer at all --- that is the entire point of
Chapter~\ref{ch:backend-architecture}'s abstraction layer. Recompiling this exact,
unmodified source file with \texttt{-DCNA\_GRAPHICS\allowbreak{}\_RENDERER=VULKAN} instead of
\texttt{OPENGLES3} produces a program that opens a Vulkan-backed window, uploads
\texttt{assets/logo.png} through the Vulkan texture path, and issues the same
\cnaclass{SpriteBatch} draw call through Vulkan command buffers instead of OpenGL calls ---
with the game-facing source unchanged. Part~\ref{part:renderers} is dedicated to what actually
differs underneath that unchanged surface, renderer by renderer.
